% Rascunho alinhado às dimensões do SBPO 2026 (ST). Substitua/mescle com o template oficial do evento quando disponível.
\documentclass[11pt,a4paper]{article}

\usepackage[a4paper,top=3.3cm,bottom=2.5cm,left=2.9cm,right=2.9cm]{geometry}
\usepackage{mathptmx}
\usepackage[T1]{fontenc}
\usepackage[utf8]{inputenc}
\usepackage[brazil]{babel}
\usepackage{microtype}
\usepackage{amsmath,amssymb}
\usepackage{booktabs}
\usepackage{graphicx}
\usepackage{pgfplots}
\pgfplotsset{compat=1.18}
\usepackage{tikz}
\usetikzlibrary{arrows.meta,positioning,calc}
\usepackage{hyperref}
\usepackage{enumitem}
\setlist{nosep}

\linespread{1}
\renewcommand{\normalsize}{\fontsize{11}{13.2}\selectfont}
\normalsize

\hypersetup{colorlinks=true,linkcolor=black,citecolor=black,urlcolor=blue}

\title{Simulação de fluxo Kanban com calendário Scrum: sinergia interpessoal, traços individuais e experimentos reprodutíveis}
\date{}

\begin{document}
\maketitle

\begin{abstract}
Equipes de software combinam \emph{sistemas puxados} ao estilo Kanban com cadências Scrum, porém modelos quantitativos explícitos costumam omitir relações interpessoais (sinergia) que afetam colaboração e repasses. Este artigo descreve um modelo operacional de simulação \textbf{em tempo discreto} (granularidade diárias) com colunas de fluxo, capacidade estocástica por membro, políticas de WIP e planejamento, e multiplicadores de sinergia para trabalho colaborativo e para repasses entre etapas, incluindo retrabalho. O motor é reprodutível via semente de PRNG. Relatamos um conjunto mínimo de cenários comparativos (A/B) com a mesma semente e \emph{backlog}, variando apenas sinergia entre pessoas, juntamente com métricas de fluxo (contagens por coluna ao longo do tempo, tempo de ciclo e \emph{throughput} por sprint). A contribuição é metodológica: um arcabouço explícito para análise de sensibilidade e ensino em pesquisa operacional, com fronteira clara entre hipóteses operacionais e validação empírica.

\end{abstract}

\noindent\textbf{Palavras-chave:} simulação a eventos discretos; fluxo de desenvolvimento de software; sinergia de equipe; reprodutibilidade experimental.

\section{Introdução}

Fluxos de desenvolvimento de software combinam filas de trabalho, políticas puxadas e cadências de planejamento~\cite{poppendieck2003lean,anderson2010kanban,reinertsen2009principles}. Em contextos ágeis, práticas inspiradas em Kanban e Scrum coexistem: um quadro visual organiza o trabalho, enquanto ritos estruturam o calendário e a transparência do progresso~\cite{schwaber2002agile,schwaber2020scrum}. Sob a ótica de pesquisa operacional, essa configuração sugere um sistema estocástico com filas, políticas de atendimento e eventos de calendário---como enfatizado pela relação $L=\lambda W$~\cite{little1961proof} e pela análise de variabilidade em sistemas com WIP limitado~\cite{hopp2008factory}; contudo, permanece difícil transmitir, em sala de aula ou em estudos exploratórios, como \emph{relações entre pares}---aqui resumidas por sinergia---interagem com a variabilidade de capacidade e com gargalos operacionais.

Simuladores educacionais comerciais, como o \emph{Kanban Simulator}~\cite{kanbanSimulator}, tornaram intuitiva a variabilidade diária e a especialização. Modelagem e simulação de processos de software são discutidas extensamente na literatura~\cite{kellner1999software,madachy2008software}, mas com foco predominante em \emph{system dynamics} ou modelos de projeto; para pesquisa e ensino avançado em PO, porém, deseja-se um modelo \textbf{enunciado em equações e parâmetros}, \textbf{reprodutível} e \textbf{separável}: deve ser possível isolar o efeito da sinergia entre dois desenvolvedores em um cartão colaborativo sem, ao mesmo tempo, alterar o \emph{backlog}, a semente aleatória ou a política de WIP.

\textbf{Objetivo.} Formalizar e apresentar um motor de simulação transparente que combina (i) fluxo Kanban diário, (ii) cadência Scrum simplificada e (iii) sinergia interpessoal na colaboração na mesma etapa e nos repasses entre etapas, permitindo experimentos A/B reprodutíveis. O texto destaca, na Seção~\ref{sec:modelo}, as \textbf{diferenças principais de abordagem} entre uma leitura \emph{sem} sinergias ($s_{ij}\equiv 0$: núcleo neutro em relação a pares) e uma leitura \emph{com} sinergias (matriz $S$ acoplada a colaboração, repasse e retrabalho).

\textbf{Contribuições.}
\begin{itemize}
  \item \textbf{Modelo operacional explícito:} tempo discreto, colunas de fluxo, capacidade por membro com componente aleatória e bônus de especialista, filas FIFO nas colunas de trabalho, e avanço entre etapas com repasses e retrabalho.
  \item \textbf{Sinergia como fator controlado:} uma matriz simétrica $s_{ij}\in[-1,1]$ entra em multiplicadores com \emph{limitação} (\emph{clamping}) na colaboração em \texttt{dev} e nos repasses entre etapas.
  \item \textbf{Calendário Scrum acoplado ao motor:} planejamento, dias úteis, revisão e retrospectiva (esta última, no protótipo, apenas como marcador de fim de sprint).
  \item \textbf{Pacote experimental mínimo:} mesma semente e \emph{backlog}, variando apenas padrões de sinergia entre pares, com exportação automática de métricas e figuras.
  \item \textbf{Artefato de software livre:} aplicação web (React/Vite) sobre o motor TypeScript, internacionalização, jogo interativo com edição de responsáveis alinhada ao motor, validação de alocação, visualizações (CFD, tempo de ciclo, relatório financeiro pedagógico) e o script \texttt{npm run paper:figures} que materializa figuras do artigo a partir do repositório.
\end{itemize}

\textbf{Organização.} A Seção~\ref{sec:relacionados} posiciona trabalhos relacionados. A Seção~\ref{sec:modelo} descreve o sistema, a notação e o algoritmo diário. A Seção~\ref{sec:ritos} liga ritos ao calendário. A Seção~\ref{sec:arquitetura} resume a arquitetura do protótipo e um fluxo visual. As Seções~\ref{sec:experimentos} e~\ref{sec:resultados} apresentam experimentos e resultados ilustrativos. A Seção~\ref{sec:discussao} discute limitações e a Seção~\ref{sec:conclusao} conclui.

\section{Trabalhos relacionados}
\label{sec:relacionados}

\paragraph{Simulação e fluxo em desenvolvimento de software.}
A literatura de simulação a eventos discretos e de teoria de filas aplica-se naturalmente a processos de conhecimento~\cite{banks2010discrete}, embora a estimação de parâmetros e a validação sejam desafios conhecidos. Em engenharia de software, métricas como throughput, tempo de ciclo e diagramas de fluxo cumulativo (CFD) tornaram-se linguagem comum em iniciativas lean/Kanban~\cite{anderson2010kanban}. Modelos pedagógicos baseados em cartões e colunas ajudam a ilustrar WIP e variabilidade, frequentemente com componente estocástica na capacidade diária.

\paragraph{Simuladores e jogos sérios.}
O \emph{Kanban Simulator}~\cite{kanbanSimulator} popularizou a dinâmica de dados por membro e bônus de especialidade na coluna correspondente, além de métricas de fluxo. Do ponto de vista de Pesquisa Operacional, tais ferramentas são valiosas como \emph{intuição}, mas nem sempre expõem todas as regras necessárias para replicação acadêmica estrita. Em paralelo, \emph{jogos sérios} e laboratórios de software têm sido usados para ensinar políticas de filas e efeitos de gargalo; nosso foco aqui é complementar: um motor declarativo, versionável e inspecionável.

\paragraph{Fatores humanos e equipes.}
Estudos em psicologia organizacional e engenharia de software discutem coesão, confiança e conflito. Nosso enquadramento não pretende substituir instrumentos psicométricos: a sinergia $s_{ij}\in[-1,1]$ é um \textbf{parâmetro} explícito do modelo, interpretável como um resumo operacional de facilitação ou atrito entre pares \emph{sob as regras do simulador}. Traços individuais funcionam como ``perfiladores'' leves que alteram coeficientes do mesmo núcleo mecânico, permitindo contrastar hipóteses sem multiplicar famílias de modelos.

\paragraph{Síntese.}
Diferentemente de abordagens puramente narrativas ou de \emph{caixas-pretas}, o arcabouço aqui descrito prioriza transparência algorítmica (implementação em TypeScript no repositório associado) e um pacote mínimo de cenários A/B para análise de sensibilidade.

\section{Descrição do sistema e modelo operacional}
\label{sec:modelo}

\subsection{Visão geral}
Consideramos um conjunto finito de \textbf{membros} $\mathcal{M}$, cada um com especialidade em $\{\text{Analista},\text{Desenvolvedor},\text{Testador}\}$, e um conjunto de \textbf{cartões} (histórias) com pontos $p\in\mathbb{N}$. O fluxo segue colunas ordenadas:
\[
\texttt{backlog}\rightarrow\texttt{ready}\rightarrow\texttt{analise}\rightarrow\texttt{dev}\rightarrow\texttt{teste}\rightarrow\texttt{deploy}.
\]
Cada cartão de pontos $p$ tem trabalho remanescente por etapa $(w^A,w^D,w^T)$ obtido por uma regra fixa de particionamento aproximadamente $30\%/50\%/20\%$ de $p$ com arredondamentos inteiros e garantia de componentes positivos (implementação: função \texttt{splitWork} no código-fonte).

\subsection{Parâmetros globais}
Além do backlog e da equipe, o modelo usa parâmetros agregados (todos auditáveis no tipo \texttt{SimulationParams}): número de dias úteis por sprint $D$, número de sprints $K$, semente $\sigma$ do PRNG, limite de WIP por coluna de trabalho $W$, limite de puxadas do backlog para \texttt{ready} no planejamento $P_{\max}$, coeficientes $\beta$ (colaboração) e $\gamma$ (handoff), extremos de \emph{clamp} para multiplicadores de colaboração e handoff, limiar $s^\star$ que controla a chance de retrabalho em handoffs desfavoráveis, e unidades de retrabalho $u_R$.

\subsection{Sinergia e traços}
Para cada par não ordenado $\{i,j\}\subset\mathcal{M}$ define-se $s_{ij}=s_{ji}\in[-1,1]$ (default $0$ se ausente). Cada membro pode carregar até uma qualidade e um defeito de um catálogo fechado; cada traço altera vetores numéricos pequenos (variação no máximo do dado, bônus de especialista, ajustes nos multiplicadores de colaboração/handoff, WIP efetivo e chance incremental de retrabalho). O efeito líquido de traços em um membro é a soma dos efeitos (função \texttt{resolveMemberModifiers}).

\begin{table}[htbp]
  \centering
  \caption{Catálogo de traços (v1) e efeitos principais no motor.}
  \label{tab:tracos}
  \begin{tabular}{@{}llp{6.8cm}@{}}
    \toprule
    Nome & Tipo & Efeito operacional resumido \\
    \midrule
    Comunicador & Qualidade & $+0{,}08$ no termo aditivo do multiplicador de handoff. \\
    Mentor & Qualidade & $+0{,}06$ no termo aditivo do multiplicador de colaboração em \texttt{dev}. \\
    Veterano & Qualidade & $+0{,}25$ no bônus de especialista (sobre a base multiplicativa). \\
    Hiperfoco & Defeito & $-1$ no máximo do dado diário (após \emph{clamp}). \\
    Evita conflito & Defeito & $-0{,}08$ no termo de colaboração em \texttt{dev}. \\
    Mentalidade de silo & Defeito & $-0{,}10$ no termo de handoff. \\
    \bottomrule
  \end{tabular}
\end{table}

\subsection{Calendário e incremento de tempo}
O tempo avança em \textbf{dias globais} $t=1,2,\dots$. A cada dia pertence a um sprint $k$ e a um índice de dia no sprint $d\in\{1,\dots,D\}$, além de um rito (Seção~\ref{sec:ritos}). O estado inclui filas por coluna e, para cada cartão em etapa de trabalho, o remanescente da etapa corrente $\textit{rem}$.

\subsection{Planejamento (puxar para \texttt{ready})}
No dia de \emph{Sprint Planning}, repetimos até $P_{\max}$ vezes: se existir cartão em \texttt{backlog} e a fila \texttt{ready} ainda tiver menos de $P_{\max}$ cartões, move-se o cartão do início de \texttt{backlog} para o fim de \texttt{ready}. Em seguida, enquanto $|\texttt{analise}|<W'$ e \texttt{ready} não estiver vazio, puxa-se cartões de \texttt{ready} para \texttt{analise} até o teto $W'$, onde $W'$ é o WIP efetivo por coluna após somar arredondamentos dos traços da equipe.

\subsection{Capacidade diária por membro}
Para cada membro $m$ no dia de trabalho, sorteia-se um inteiro uniforme
\[
X_m \sim \mathcal{U}\{1,\dots,D^{\max}_m\},\qquad
D^{\max}_m = \mathrm{clip}\big(6+\Delta^{\text{dado}}_m,\,2,\,8\big),
\]
em que $\Delta^{\text{dado}}_m$ agrega traços. Seja $h_m\in\{\texttt{analise},\texttt{dev},\texttt{teste}\}$ a coluna da especialidade de $m$. Se existir trabalho remanescente positivo em $h_m$, define-se capacidade bruta
\[
C_m = \min\left\{12,\left\lfloor X_m \cdot \rho_m + \tfrac{1}{2}\right\rfloor\right\},
\quad
\rho_m = \mathrm{clip}\big(2+\Delta^{\text{esp}}_m,\,1.25,\,3.5\big),
\]
caso contrário $C_m=X_m$. O valor $C_m$ é a \textbf{capacidade bruta} aplicável às colunas de trabalho.

\subsection{Política de alocação da capacidade}
Para cada membro $m$ nesta ordem, inicializa-se $L\leftarrow C_m$. Primeiro, se sua especialidade $h_m$ possui cartões com $\textit{rem}>0$, aplica-se o procedimento \textsc{Spend} em $h_m$ (abaixo). Depois, para cada coluna de trabalho na ordem $(\texttt{analise},\texttt{dev},\texttt{teste})$, enquanto $L>0$, aplica-se \textsc{Spend}.

\paragraph{Procedimento \textsc{Spend} em uma coluna $h$.}
Percorre-se a fila de $h$ na ordem FIFO. Para cada cartão com $\textit{rem}>0$, define-se multiplicador de colaboração
\[
M^{\text{col}} =
\begin{cases}
\mathrm{clip}\big(1+\beta\, s_{ab} + T^{\text{col}}_{ab},\,E^{\min}_{\text{col}},\,E^{\max}_{\text{col}}\big), & \text{se colaborativo com assignees }a,b,\\
1, & \text{caso contrário},
\end{cases}
\]
onde $T^{\text{col}}_{ab}$ agrega termos de traço dos dois assignees. Na coluna \texttt{dev} apenas, o trabalho efetivo consumido por unidade bruta é multiplicado por $M^{\text{col}}$; nas demais, $M^{\text{col}}=1$. Dado orçamento bruto $L$, para um cartão com remanescente $\textit{rem}$ o máximo de unidades brutas aplicáveis é $\textit{rem}/M^{\text{col}}$ em \texttt{dev} (e $\textit{rem}$ nas outras). Consome-se $\delta=\min\{L,\textit{rem}/M^{\text{col}}\}$ (com convenção acima), reduzindo $\textit{rem}\leftarrow \textit{rem}-\delta M^{\text{col}}$ em \texttt{dev} (e $\textit{rem}-\delta$ nas outras). Se $\textit{rem}$ zera, tenta-se avançar o cartão.

\subsection{Avanço entre colunas, handoff e retrabalho}
Ao avançar de uma coluna de trabalho $h$ para a próxima, sejam $o$ e $i$ os membros ``doador'' e ``receptor'' inferidos pelas especialidades das colunas (com regras de fallback para assignees e equipe mínima). Define-se
\[
M^{\text{hand}} = \mathrm{clip}\big(1+\gamma\, s_{oi} + T^{\text{hand}}_{oi},\,H^{\min},\,H^{\max}\big),
\]
com $T^{\text{hand}}$ somando termos de traço. Com probabilidade
\[
\pi_R = \mathrm{clip}\Big(\mathbb{I}\{s_{oi}<s^\star\}\cdot\big(0.06+(s^\star-s_{oi})\,0.22\big) + \Delta^{\text{rw}}_o+\Delta^{\text{rw}}_i,\,0,\,0.55\Big),
\]
sorteia-se retrabalho; se ocorrer, acrescentam-se $u_R$ unidades ao estágio de destino (e registra-se nota explicativa). Ao entrar em \texttt{dev} ou \texttt{teste} a partir de um handoff entre especialidades, o remanescente inicial do estágio é majorado por arredondamentos que incorporam $M^{\text{hand}}$ (detalhe implementado no código-fonte).

\subsection{Métricas}
Do registro diário de contagens por coluna (estrutura \texttt{DayLog} no código) constrói-se uma série temporal por coluna, útil para visualizações no estilo de CFD e para inspeção de gargalos. Para cada cartão concluído, mede-se tempo de ciclo (dias globais entre entrada em \texttt{ready} e chegada a \texttt{deploy}). O throughput por sprint agrega conclusões pelo sprint associado ao dia global de \texttt{deploy}.

\subsection{Pseudocódigo do dia de trabalho}
A listagem abaixo condensa a política de alocação descrita nas subseções anteriores; detalhes de handoff, retrabalho e mensagens de diagnóstico permanecem no código-fonte.

\begin{figure}[htbp]
  \centering
  \begin{minipage}{0.92\linewidth}
  \small
  \begin{verbatim}
Para cada membro m (ordem fixa da equipe):
  sortear X_m e calcular C_m (aplicar multiplicador de especialista se houver
                             trabalho remanescente na coluna da especialidade)
  L <- C_m
  Gastar(h_m, L)                    // prioriza a coluna da especialidade
  para h em (analise, dev, teste):
    Gastar(h, L)

Gastar(h, L):
  para cada cartão na fila h (FIFO, cópia do estado atual):
    se rem <= 0: continuar
    se h == dev e cartão colaborativo:
       M <- M_col(sinergia, traços dos responsáveis)
       consumir delta <- min(L, rem/M) em unidades brutas
       rem <- rem - delta * M
    senão:
       consumir delta <- min(L, rem)
       rem <- rem - delta
    se rem == 0: tentar avançar cartão (handoff, possível retrabalho)
  \end{verbatim}
  \end{minipage}
  \caption{Visão compacta do passo de trabalho diário (processamento sequencial por membro; cada um consome sua capacidade bruta antes do próximo).}
  \label{fig:pseudo}
\end{figure}

\subsection{Implementação e reprodutibilidade}
O motor é implementado em TypeScript puro (sem dependência de UI) no diretório \texttt{src/simulation/}, expondo \texttt{runSimulation(config)} com PRNG determinístico por $\sigma$ (\texttt{rng.ts}). O arquivo de cenários experimentais e parâmetros fixos espelha-se em \texttt{scripts/paperScenarios.ts}, e figuras/tabelas podem ser regeneradas com \texttt{npm run paper:figures}. Isso facilita auditoria por revisores e replicação em aulas ou laboratórios.

\section{Ritos Scrum no calendário}
\label{sec:ritos}

Para cada sprint $k=1,\dots,K$, indexamos dias $d=1,\dots,D$ e um dia adicional formal de \emph{retrospectiva} sem alteração mecânica do estado no protótipo atual (apenas registro em log).

\begin{table}[htbp]
  \centering
  \caption{Mapeamento simples entre dia no sprint e modo de processamento.}
  \label{tab:ritos}
  \begin{tabular}{@{}llp{7.2cm}@{}}
    \toprule
    Dia $d$ & Rito / modo & Efeito no motor (v1) \\
    \midrule
    $1$ & \emph{Sprint Planning} & Puxa cartões de \texttt{backlog} para \texttt{analise} até $P_{\max}$ (e WIP); em seguida, se $P_{\max}>0$, continua a preencher \texttt{analise} a partir do \texttt{backlog} até ao WIP. \\
    $2,\dots,D-1$ & \emph{Daily Scrum} & Dia de trabalho: aplica capacidade, alocação e avanços. \\
    $D$ & \emph{Sprint Review} & Também dia de trabalho (mesma mecânica), com marcação explícita no log. \\
    $D{+}1$ & \emph{Retrospectiva} & Sem efeito mecânico; encerra sprint e reinicia $d\leftarrow 1$. \\
    \bottomrule
  \end{tabular}
\end{table}

A separação pedagógica entre ``marcador'' de Daily e o processamento de trabalho facilita narrativas de ensino: o calendário Scrum aparece como \textbf{ritmo} sobreposto ao \textbf{fluxo} Kanban~\cite{schwaber2002agile}, sem exigir modelagem de horas de reunião na primeira versão. Em termos de modelagem OR, os ritos funcionam como \textbf{modos} do mesmo relógio de simulação: mudam quais transições são habilitadas no início do dia (puxar \emph{vs.} apenas processar), preservando um único estado de sistema entre dias.

\paragraph{Eventos aleatórios de capacidade (Daily e Review).}
Nos dias de \emph{Daily Scrum} e \emph{Sprint Review}, o motor aplica uma probabilidade global configurável (\texttt{dailyRandomEventChance}): com essa probabilidade, sorteia-se \emph{no máximo} um evento de um catálogo ponderado (doença, férias, meio período, reunião de equipa, incidente de infraestrutura, ``bom clima''), cada qual com multiplicador ou anulação de capacidade efetiva para membros afetados. Os acontecimentos aparecem no \texttt{DayLog} e são listados na interface para consulta pedagógica.

\section{Arquitetura do protótipo e fluxo visual}
\label{sec:arquitetura}

Do ponto de vista de engenharia de software, o protótipo separa um \textbf{núcleo de simulação} (TypeScript puro, sem dependência de \emph{framework}) de uma camada de interface que pode evoluir para um fluxo em duas fases (configuração do cenário e execução da simulação, no espírito de \emph{setup} e \emph{partida}) sem alterar o contrato do motor. Na versão atual do repositório, o núcleo concentra tipos, catálogo de traços, resolução de modificadores, PRNG com semente, sinergia e o motor diário (\texttt{runSimulation}), enquanto rotinas auxiliares geram séries para CFD e tempos de ciclo.

\begin{figure}[htbp]
  \centering
  \begin{tikzpicture}[
    col/.style={draw, rounded corners=2pt, minimum width=1.05cm, minimum height=0.9cm, align=center, font=\scriptsize},
    arr/.style={-{Latex[length=2mm]}, thick},
    note/.style={font=\scriptsize, align=center},
  ]
    \node[col] (bl) {Backlog};
    \node[col, right=0.35cm of bl] (rd) {Ready};
    \node[col, right=0.35cm of rd] (an) {Análise};
    \node[col, right=0.35cm of an] (dv) {Dev};
    \node[col, right=0.35cm of dv] (ts) {Teste};
    \node[col, right=0.35cm of ts] (dp) {Deploy};
    \draw[arr] (bl) -- (rd);
    \draw[arr] (rd) -- (an);
    \draw[arr] (an) -- (dv);
    \draw[arr] (dv) -- (ts);
    \draw[arr] (ts) -- (dp);
    \path (bl.north west) ++(-0.15,0.55) coordinate (topL);
    \path (dp.north east) ++(0.15,0.55) coordinate (topR);
    \draw[dashed, gray] (topL) rectangle ($(dp.south east)+(0.15,-0.35)$);
    \node[note, anchor=south] at ($(topL)!0.5!(topR)$) {Quadro Kanban (colunas de fluxo)};
    \node[note, anchor=north west] at ($(bl.south west)+(-0.1,-0.85)$)
      {\textbf{Ritmo Scrum (v1):} dia $1$ --- planejamento (puxar para \emph{Ready} e Análise);\\
       dias $2,\ldots,D-1$ --- trabalho (\emph{Daily});\\
       dia $D$ --- trabalho e revisão do sprint (marcador);\\
       dia formal de retrospectiva sem efeito mecânico.};
  \end{tikzpicture}
  \caption{Fluxo principal de cartões e relação conceitual com o calendário de ritos.}
  \label{fig:fluxoQuadro}
\end{figure}

A Figura~\ref{fig:fluxoQuadro} resume a leitura pedagógica que queremos preservar: o Kanban organiza \emph{o que} está em cada estágio; o Scrum organiza \emph{quando} certas ações de cadência ocorrem. O motor materializa essa combinação ao alternar transições de planejamento e passos de trabalho conforme a Tabela~\ref{tab:ritos}.

\section{Desenho experimental}
\label{sec:experimentos}

\subsection{Objetivo e hipóteses operacionais}
Ilustrar \textbf{análise de sensibilidade} sob um modelo transparente: mantemos fixos o backlog inicial, a semente do PRNG, a política de WIP e os parâmetros globais (Tabela~\ref{tab:param}), variando apenas (i) sinergia em um par colaborativo em \texttt{dev}, (ii) sinergia negativa no handoff entre analista e desenvolvedor principal, e (iii) um defeito individual (\emph{Hiperfoco}) no desenvolvedor principal.

As hipóteses \emph{operacionais} associadas a este desenho são modestas e \textbf{dependentes dos parâmetros}: esperamos que (H1) sinergia positiva no par colaborativo reduza, em média, o tempo de ciclo dos cartões que se beneficiam de maior produtividade efetiva em \texttt{dev}; (H2) sinergia negativa no handoff da analista para o desenvolvedor principal aumente atrasos via combinação de multiplicador de handoff e maior probabilidade de retrabalho; (H3) o defeito \emph{Hiperfoco} reduza throughput médio ao restringir o suporte do dado diário, com impacto variável na média de ciclo conforme a realização estocástica e a contenção do sistema. A Seção~\ref{sec:resultados} confronta essas expectativas com números gerados automaticamente.

\subsection{Equipe e backlog}
Utilizamos quatro membros: uma analista, dois desenvolvedores (para permitir cartões colaborativos) e um testador. O backlog possui sete cartões com pontuações heterogêneas; dois cartões são colaborativos com assignees nos dois desenvolvedores. Essa composição força, de forma controlada, o uso do mecanismo de colaboração em \texttt{dev} sem alterar o tamanho total do trabalho entre cenários.

\subsection{Cenários}
\begin{itemize}
  \item \textbf{A (linha de base):} sinergias nulas e sem traços.
  \item \textbf{B (colaboração):} $s=0{,}88$ entre os dois desenvolvedores; demais sinergias nulas.
  \item \textbf{C (handoff):} $s=-0{,}82$ entre analista e desenvolvedor principal; demais sinergias nulas.
  \item \textbf{D (traço):} como A, porém o desenvolvedor principal recebe o defeito \emph{Hiperfoco}.
\end{itemize}

\subsection{Parâmetros numéricos fixos}
\begin{table}[htbp]
  \centering
  \caption{Parâmetros globais usados nos experimentos (gerados a partir do código).}
  \label{tab:param}
  % Gerado por scripts/generate-paper-figures.ts — não editar à mão
\begin{tabular}{@{}lr@{}}
  \toprule
  Parâmetro & Valor \\
  \midrule
  \texttt{daysPerSprint} & 10 \\
  \texttt{numSprints} & 10 \\
  \texttt{seed} & 424242 \\
  \texttt{wipPerColumn} & 3 \\
  \texttt{planningPullMax} & 6 \\
  \texttt{synergyBeta} & 0.38 \\
  \texttt{synergyGamma} & 0.42 \\
  \texttt{collabEffMin} / \texttt{collabEffMax} & 0.72 / 1.38 \\
  \texttt{handoffEffMin} / \texttt{handoffEffMax} & 0.68 / 1.28 \\
  \texttt{handoffReworkSynergyThreshold} & -0.12 \\
  \texttt{reworkUnits} & 2 \\
  \bottomrule
\end{tabular}

\end{table}

\subsection{Métricas reportadas e reprodutibilidade}
Para cada cenário, registramos: (i) série diária de contagens por coluna (base para visualizações estilo CFD); (ii) tempos de ciclo por cartão concluído (de \texttt{ready} a \texttt{deploy}), com média e mediana por cenário; (iii) throughput por sprint, atribuindo cada conclusão ao sprint do dia global correspondente. Todos os arquivos em \texttt{paper/data/} e \texttt{paper/figures/} podem ser regenerados a partir do mesmo commit do código, o que reduz risco de descompasso entre texto e implementação.

\section{Resultados ilustrativos}
\label{sec:resultados}

Todos os números desta secção são \textbf{gerados automaticamente} por \texttt{npm run paper:figures} a partir do motor \texttt{runSimulation} do repositório. As figuras usam \texttt{pgfplots} lendo ficheiros CSV exportados pelo mesmo script, de modo que pequenos ajustes de parâmetros podem refletir-se rapidamente no PDF.

% Gerado por scripts/generate-paper-figures.ts
\newcommand{\StatCompletedA}{7}
\newcommand{\StatMeanCtA}{2.71}
\newcommand{\StatMedianCtA}{3.00}
\newcommand{\StatCompletedB}{7}
\newcommand{\StatMeanCtB}{2.57}
\newcommand{\StatMedianCtB}{2.00}
\newcommand{\StatCompletedC}{7}
\newcommand{\StatMeanCtC}{3.00}
\newcommand{\StatMedianCtC}{3.00}
\newcommand{\StatCompletedD}{7}
\newcommand{\StatMeanCtD}{2.71}
\newcommand{\StatMedianCtD}{3.00}


A Figura~\ref{fig:cfdA} mostra a dinâmica de filas no cenário A: como o estoque em \texttt{backlog} se relaciona com o trabalho em progresso nas colunas de especialidade. As Figuras~\ref{fig:bar} e~\ref{fig:tp} resumem comparações entre cenários: a primeira destaca o tempo de ciclo médio (a mediana complementa a Tabela~\ref{tab:resumoCenarios}), enquanto a segunda mostra diferenças de \emph{throughput} por sprint quando o padrão de sinergia entre pares muda.

\begin{figure}[htbp]
  \centering
  % Generated by scripts/generate-paper-figures.ts
\begin{tikzpicture}
\begin{axis}[
  width=0.95\linewidth,
  height=5.4cm,
  xmin=1,
  xlabel={Dia global},
  ylabel={Cartões por coluna},
  legend columns=3,
  legend style={font=\scriptsize,at={(0.5,1.02)},anchor=south},
]
\addplot[thick,const plot,color=black!70] table[col sep=comma,x=day,y=backlog]{data/cfd_A_baseline.csv};
\addplot[thick,const plot,color=violet!80] table[col sep=comma,x=day,y=analise]{data/cfd_A_baseline.csv};
\addplot[thick,const plot,color=teal!80] table[col sep=comma,x=day,y=dev]{data/cfd_A_baseline.csv};
\addplot[thick,const plot,color=orange!90] table[col sep=comma,x=day,y=teste]{data/cfd_A_baseline.csv};
\addplot[thick,const plot,color=red!80] table[col sep=comma,x=day,y=deploy]{data/cfd_A_baseline.csv};
\legend{Backlog,Análise,Dev,Testes,Deploy}
\end{axis}
\end{tikzpicture}

  \caption{Contagens de cartões por coluna ao longo dos dias globais no cenário A (referência).}
  \label{fig:cfdA}
\end{figure}

\begin{figure}[htbp]
  \centering
  % Generated by scripts/generate-paper-figures.ts
\begin{tikzpicture}
\begin{axis}[
  width=0.62\linewidth,
  height=5cm,
  ybar,
  ymin=0,
  symbolic x coords={A,B,C,D},
  xtick=data,
  xticklabel style={align=center},
  ylabel={Ciclo médio (dias)},
  nodes near coords,
  nodes near coords align={vertical},
  bar width=16pt,
]
\addplot[fill=blue!55] coordinates {(A,1.86) (B,1.71) (C,2.43) (D,1.86)};
\end{axis}
\end{tikzpicture}

  \caption{Tempo de ciclo médio (dias) por cenário, considerando apenas cartões concluídos na janela simulada.}
  \label{fig:bar}
\end{figure}

\begin{figure}[htbp]
  \centering
  % Generated by scripts/generate-paper-figures.ts
\begin{tikzpicture}
\begin{axis}[
  width=0.95\linewidth,
  height=5.2cm,
  xlabel={Sprint},
  ylabel={Cartões concluídos},
  xmin=1,
  xmax=10,
  ymin=0,
  legend style={font=\scriptsize,at={(0.5,1.02)},anchor=south},
  legend columns=4,
]
\addplot[thick,mark=*,color=black!70] table[col sep=comma,x=sprint,y=A]{data/throughput_by_sprint.csv};
\addplot[thick,mark=square*,color=blue!70] table[col sep=comma,x=sprint,y=B]{data/throughput_by_sprint.csv};
\addplot[thick,mark=triangle*,color=red!75] table[col sep=comma,x=sprint,y=C]{data/throughput_by_sprint.csv};
\addplot[thick,mark=diamond*,color=teal!80] table[col sep=comma,x=sprint,y=D]{data/throughput_by_sprint.csv};
\legend{A,B,C,D}
\end{axis}
\end{tikzpicture}

  \caption{\emph{Throughput} por sprint (cartões que chegam a \texttt{deploy} nessa sprint) para os quatro cenários.}
  \label{fig:tp}
\end{figure}

\begin{table}[htbp]
  \centering
  \caption{Resumo dos quatro cenários (ciclo médio e mediano em dias; valores gerados em \texttt{summary\_stats.tex}).}
  \label{tab:resumoCenarios}
  \begin{tabular}{@{}lccc@{}}
    \toprule
    Cenário & Concluídos & Ciclo médio & Ciclo mediano \\
    \midrule
    A & \StatCompletedA & \StatMeanCtA & \StatMedianCtA \\
    B & \StatCompletedB & \StatMeanCtB & \StatMedianCtB \\
    C & \StatCompletedC & \StatMeanCtC & \StatMedianCtC \\
    D & \StatCompletedD & \StatMeanCtD & \StatMedianCtD \\
    \bottomrule
  \end{tabular}
\end{table}

\paragraph{Leitura rápida.}
No pacote aqui reportado, o cenário B tende a melhorar o tempo de ciclo médio em relação a A quando a colaboração em \texttt{dev} é amplificada, enquanto o cenário C piora a média ao combinar um \emph{repasse} desfavorável com maior propensão a retrabalho. O cenário D combina os dois efeitos de $S$ e permite discutir interação entre canais de sinergia para a mesma semente e \emph{backlog}.

Ficheiros adicionais em \texttt{paper/figures/} (SVG) e \texttt{paper/data/} (CSV/JSON) permitem reproduzir gráficos fora do \LaTeX{}.

\section{Discussão e limitações}
\label{sec:discussao}

\paragraph{Interpretação.}
Os resultados da Seção~\ref{sec:resultados} não devem ser lidos como ``verdades'' sobre equipes reais: mostram \textbf{comportamento sensível} às regras explícitas do motor sob um desenho experimental mínimo. Recorde-se a distinção da Seção~\ref{sec:sinergia-neutral-vs}: relativamente à visão com $S\equiv 0$, introduzir sinergias não ``muda o jogo'' de propósito (continuam FIFO, WIP e sorteios de capacidade), mas \textbf{relocaliza} parte da variabilidade de fluxo em pares explícitos e nos mecanismos de colaboração/repasse. Em particular, a sinergia negativa de repasse (cenário C) combina o multiplicador de repasse com maior propensão a retrabalho, enquanto a sinergia positiva no par colaborativo (cenário B) amplifica o trabalho efetivo apenas na coluna \texttt{dev} para cartões colaborativos. O cenário D ilustra a combinação dos dois canais de $S$, reforçando a necessidade de desenhos experimentais mais amplos quando o objetivo é estimar magnitude com precisão.

\paragraph{Reprodutibilidade e transparência.}
Manter o motor em TypeScript puro e versionar em conjunto (i) parâmetros, (ii) cenários, (iii) scripts de exportação do artigo e (iv) a aplicação web que consome o mesmo pacote \texttt{src/simulation} reduz ambiguidade sobre ``o que foi simulado'', alinhando-se às práticas recomendadas por Wilson et al.~\cite{wilson2017good} e às regras de Sandve et al.~\cite{sandve2013ten} para pesquisa computacional reprodutível. Trabalhos futuros devem publicar \emph{releases} com \texttt{package-lock.json} (ou equivalente) e um identificador de \emph{commit} ao reproduzir números em artigos ou materiais didáticos.

\paragraph{Limitações.}
O modelo é redutivo: sem aprendizado, sem filas paralelas de defeitos, sem mudanças de prioridade no meio da sprint, e a retrospectiva não altera parâmetros. A sinergia $s_{ij}$ é um escalar-resumo por par.

\paragraph{Validação.}
Não reportamos validação empírica com dados de projetos reais; esse é um passo natural seguinte (calibração de $\beta$ e $\gamma$), como discutido por Kellner et al.~\cite{kellner1999software} no contexto de simulação de processos de software.

\section{Conclusão}
\label{sec:conclusao}

Apresentamos um simulador de fluxo Kanban com cadência Scrum, sinergia par a par e traços individuais, com foco em \textbf{transparência algorítmica} e \textbf{reprodutibilidade} via semente de PRNG. O texto combina formalização (Seção~\ref{sec:modelo}), calendário de ritos (Seção~\ref{sec:ritos}), arquitetura do protótipo (Seção~\ref{sec:arquitetura}) e um pacote mínimo de cenários A/B com mesma carga inicial e mesmos parâmetros globais, permitindo discutir sensibilidade de métricas de fluxo a fatores interpessoais operacionalizados de forma explícita.

Em síntese, a contribuição central não é ``prever'' um time real, mas disponibilizar um \textbf{laboratório computacional} com fronteira clara entre hipótese operacional e evidência empírica. Como extensões, citamos: calibração com dados de projetos, catálogo de traços mais rico (com testes de face validity), políticas alternativas de priorização, e integração com protocolos de experimentação com equipes humanas jogando o simulador.


\bibliographystyle{plain}
\bibliography{bibliografia}

\end{document}
